\section{The Three Interpretations Compared}
\label{sec:comparison}

All three interpretations read the parable as a call that is wholly God's gift
and that can be refused or wasted. The first and the third end in a fullness,
the first in the feast completed at the resurrection and the third in desire
filled in glory; the second ends at the king's inspection and in a vision
measured by charity. They differ in the question they put to the parable, and
therefore in the verses on which they rest, in what they make of the rest of
the Mass, and in what they ask of the hearer.

\begin{comparisontable}{Question}{The king's feast for all peoples}{The garment the called must wear}{Why the invited would not come}
Which verses carry it? & Mt 22:2, 4, 8--10 with Is 25:6--9, the verses of the
\work{Ordo}'s two correlated titles; in the longer form also v.~11. & Mt 22:11--14, in the longer form only. &
Mt 22:3, 5--6, in both forms.\\
What is the feast? & The Son's wedding with the Church (Gregory, Jerome),
spread now (Gregory; Augustine) and completed at the resurrection (Hilary;
Aquinas). & Taken as given; the question is how the guest is fit for it. & The
mystery of the Incarnation, which the preoccupied will not weigh (Gregory).\\
Who is excluded, and why? & Those who will not come; several Fathers identify
the first invited with a people, a reading presented under the Council's
bound. & The guest without the garment, for lack of charity (Augustine,
Gregory) or of the garment received and not kept (Chrysostom, Hilary). & Those
held by farm and trade, by gain and ambition (Gregory, Hilary, Aquinas).\\
What does the Collect say? & Grace goes before that we may be called and
follows that we may be glorified (Aquinas quoting Augustine): the call and the
feast. & Grace goes before and follows and keeps us intent on good works:
heard as the prayer for the garment. & Intent on good works,
not on farm and trade.\\
What do the Communion antiphons say? & First: the rich lack the Living Bread
(Augustine). Second: the King seen as he is (Aquinas; Augustine), and in the
longer form the King who comes in to see the guests. & First: the
heart's riches of virtue (Augustine). Second: the hope that purifies; the more
charity, the fuller the vision (Augustine; Aquinas). & First: the rich hunger,
the seekers lack nothing (Augustine; Bellarmine). Second: desire emptied of
the world to be filled (Augustine).\\
Principal witnesses & Gregory, Jerome, Hilary, Irenaeus, Aquinas. &
Augustine, Gregory, Chrysostom, Irenaeus, Aquinas, Jerome, Hilary, who name
the garment differently. & Gregory, Chrysostom, Aquinas, Augustine; Hilary, Jerome and
Bellarmine supporting.\\
Under the shorter Gospel & Holds without v.~11: the two feasts are joined by
the resurrection alone, without the judgment between them. & Holds only in
part: its Gospel centre is in vv.~11--14. & Holds.\\
\end{comparisontable}

The first interpretation rests on the Lectionary's one official correlation,
the feast of Isaiah with the wedding of the Gospel, and on witnesses who
expound one or both of those readings; Aquinas's joining of Isaiah's feast to
the king's dinner belongs to it. The second rests on the Gospel's closing
verses, read only in the longer form, and on the Collect; Isaiah enters it
only by the editor's verbal link of death swallowed up, which no witness of
the garment draws. The third rests on the Gospel's fifth verse and reads the
rest of the formulary from the second reading and the first Communion
antiphon, which the books join to the Gospel by no relation at all and no
witness joins to it.

They answer different questions about different verses, and so they reach
different consequences. The first asks what God gives, and answers with the
Church's present table and the resurrection; its consequence is hope and the
universal reach of the call, and it must carry the parable's hardest verse,
the burnt city, under the Council's bound. The second asks what those who
have been gathered must have, and answers with a garment given and kept,
which Augustine and Gregory name charity, and with the grace that follows; its
consequence is humility and the works of love. The third asks
what keeps the invited away, and answers with goods that cannot fill; its
consequence is detachment and the right use of plenty and want. No witness
cited denies another interpretation's claim, and Gregory and Aquinas each
contribute to all three from different sentences: Gregory from his account of
the wedding, of the garment and of the farm, Aquinas from his Isaiah, his
garment of Christ and his refusal of every temporal excuse.

The interpretations also hear the same texts differently, and the
differences are instructive. The Collect's grace that goes before and follows
is, in the first, the call and the glory of the feast; in the second, the
grace that clothes the guest; in the third, the attention turned from farm and
trade. The second Communion antiphon's vision is, in the first, the King seen
as he is; in the second, the vision measured by charity; in the third, the
desire stretched by waiting. The Prayer after Communion's share in the divine
nature is, in the first, the wedding of the Church shared at the table; in
the second, the life of grace that the garment is; in the third, the flight
from the concupiscence of the world that the verse of 2~Peter names.

Three real disputes among the witnesses cross these lines. What the wedding garment is divides Augustine and Gregory, who
name charity and deny that it is baptism or faith, from Hilary, whose garment
is the glory of the Spirit received at the confession of the good questioning,
on the reading of his phrase as an allusion to baptism. This touches the
second interpretation at its centre. Whether the wedding is the present Church
or the last banquet is a difference of emphasis between Gregory, who reads the
parable of the present Church, and Hilary, who reads it of the glory of the
resurrection, with Augustine's two feasts and Aquinas's four expositions
between them; it touches the first interpretation, which joins the two by the
resurrection and, where the longer form is read, by the judgment, and does not
identify them. Whether the burnt city is Jerusalem under
Titus or the eternal fire divides Chrysostom from Gregory and Hilary, with
Jerome and Aquinas offering both; it touches the first and the third, and
under either answer the identification of the first invited remains the
reading of the Fathers who make it.

The shorter form of the Gospel changes the balance. Where it is read, the
third interpretation holds whole. The first loses the king's coming in to see
the guests: it keeps the feast spread now and the feast completed in the
resurrection, joined by the resurrection alone and without the judgment
between them, and the second Communion antiphon answers the King of the
parable's opening rather than his inspection. The second keeps only the
gathering of bad and good, without the garment that is its answer. The two
Communion antiphons change it too: each interpretation hears both, but the
first antiphon carries most weight in the third interpretation and the second
in the first.
